\documentclass[letter,11pt]{article}
\usepackage[utf8]{inputenc}

\usepackage{TLCresume}

%====================
% CONTACT INFORMATION
%====================
\def\name{Konstantin Volodin}
\def\phone{+1-438-370-2046}
\def\city{Montréal, QC, Canada}
\def\email{volodin.kostia@gmail.com}
\def\role{ML Engineering}

%====================
% HEADER
%====================
%====================
% Header: Contact
%====================

\RequirePackage{fancyhdr}
\fancypagestyle{fancy}{%
\fancyhf{}
\lhead{
		\href{mailto:\email}{\email} \\ 
		\phone \\ % PHONE    
		\city  }% city
	\chead{%
	    \centering {\Huge \skills \name \vspace{.25em}} \\% feel free to adjust vspace to 0 
     
	    {\color{highlight} \textbf{\Large{\role}}}}%
     
	   \rhead{\href{https://konstantin-volodin.github.io/}{konstantin-volodin.github.io}\\ % Portfolio
	    \href{https://github.com/Konstantin-Volodin/}{github.com/Konstantin-Volodin} \\% Github
				\href{https://www.linkedin.com/in/konstantin-volodin/}{linkedin.com/konstantin-volodin}}
	    
\renewcommand{\headrulewidth}{0pt}% 2pt header rule
\renewcommand{\headrule}{\hbox to\headwidth{%
	\color{highlight}\leaders\hrule height \headrulewidth\hfill}}
}
\pagestyle{fancy}

\setlength{\headheight}{80pt}
\setlength{\headsep}{5pt}


\begin{document}

\setlength{\parindent}{0pt}
\setlength{\parskip}{0.8em}

\section{}
\vspace{1.5em}

Hi,

I'm interested in the Senior ML Engineer, ML Ops Framework role. For the past three years, I've been building and running the highest-volume ML platform at Immigration, Refugees and Citizenship Canada - about 150,000 inference requests a week flowing through PySpark and PyArrow pipelines over Parquet, with model versioning, cross-run comparison workflows, and a deployment pipeline I manage end-to-end on AWS.

What excites me about your team is the mission of enabling ML developers to move faster. That's what I've been doing - designing a model serving architecture with 40+ SQL views, building model comparison tooling against versioned S3 artifacts, and cutting our core pipeline runtime by 60\%. I want to do that work at a larger scale, for a team pushing the boundaries of what ML systems can do.

I don't have Ray or L4 experience - I want to be upfront about that. What I do bring is production ML platform ownership, strong AWS and IaC fundamentals, and a track record of learning infrastructure fast. I taught myself the full containerized deployment stack (Docker, ECR, CDK, Fargate) on my own when the knowledge wasn't available on my team. I'm comfortable being dropped into complex systems and figuring out how they work.

I'm based in Montreal and would welcome the chance to discuss how I can contribute to the team.

\vspace{1.5em}
Konstantin

\end{document}
